\subsection{
    Segurança de Arquiteturas Modernas
}\label{sec:seguranca_arquiteturas_modernas}

A evolução dos sistemas web para arquiteturas modernas, como microsserviços e aplicações baseadas em nuvem, transformou significativamente a superfície de ataque. Com a fragmentação de aplicações monolíticas em Serviços Web distribuídos, a comunicação inter-processos e a exposição de APIs tornaram-se o núcleo do funcionamento e, consequentemente, da segurança destas aplicações.

\subsubsection{
    APIs REST — OWASP API Top 10, BOLA, \textit{rate limiting}
}\label{sec:apis_rest}

APIs REST são predominantes em aplicações modernas e concentram riscos relevantes. Segundo o OWASP API Top 10 (2023), a vulnerabilidade de maior destaque é \textit{Broken Object Level Authorization} (BOLA), que ocorre quando o backend expõe identificadores de objetos sem validar se o usuário autenticado possui permissão para acessá-los. Esse cenário permite leitura, alteração ou exclusão indevida de dados de terceiros.

Outro ponto crítico é a ausência de \textit{rate limiting}, associada ao \textit{Unrestricted Resource Consumption} (Top 4 em 2023). Sem limites de volume e frequência, requisições podem esgotar recursos da aplicação e de serviços dependentes, degradando disponibilidade e ampliando a superfície de abuso.

\subsubsection{
    APIs GraphQL — introspecção, \textit{query depth}, DoS
}\label{sec:apis_graphql}

GraphQL, embora amplamente adotado, introduz riscos específicos de segurança. Como o cliente define a complexidade da requisição em um único \textit{endpoint} (geralmente \texttt{/graphql}), consultas sintaticamente válidas podem impor alto custo computacional e causar DoS. Um vetor recorrente é a \textit{query depth}, isto é, o nível de aninhamento dos campos. A mitigação inclui validação estrutural da consulta e análise da \textit{Abstract Syntax Tree} (AST) antes da execução.

A introspecção, recurso nativo do GraphQL, permite descobrir toda a estrutura da API. Em produção, sua exposição facilita ataques mais precisos, incluindo BOLA e consultas profundas para exaustão de recursos. Por isso, recomenda-se desativar a introspecção em ambientes produtivos.

\subsubsection{
    Microserviços — mTLS, API \textit{Gateway}, gestão de segredos
}\label{sec:microservicos}

O artigo "Zero Trust Security Architecture in Microservices via mTLS" (IJARCSE, 2025) propõe eliminar a confiança implícita em redes distribuídas. Nesse modelo, o mTLS (\textit{Mutual Transport Layer Security}) é o pilar da arquitetura Zero Trust, garantindo criptografia e autenticação mútua entre serviços.

\textit{Sidecars} (como Envoy) validam certificados a cada conexão, enquanto o API Gateway centraliza o controle de acesso externo e valida tokens JWT com o provedor de identidade (por exemplo, Keycloak). Após validação, o tráfego segue para os microsserviços por canais mTLS. A gestão de segredos é reforçada por rotação frequente de certificados X.509 (como no Citadel/Istio) e uso de tokens de curta duração, reduzindo a janela de exploração em caso de comprometimento.

\subsubsection{
    \textit{Cloud-native} — IaC \textit{scanning}, \textit{Kubernetes}, S3
}\label{sec:cloud_native}

Aplicações \textit{cloud-native} dependem de infraestrutura dinâmica e altamente configurável, o que amplia riscos de \textit{misconfiguration}. Um cenário recorrente é a exposição pública indevida de recursos em provedores de nuvem, como \textit{buckets} S3, possibilitando acesso não autorizado a dados sensíveis. A mitigação exige \textit{hardening} contínuo e revisão sistemática de permissões (ACLs e políticas de acesso).

Esse ecossistema também depende de IaC e \textit{pipelines} de CI/CD, alinhando-se ao risco de integridade de software destacado pela OWASP (A08:2021). Sem controles robustos, \textit{pipelines} podem introduzir código malicioso ou acessos indevidos. Defesas recomendadas incluem governança de acesso ao CI/CD, verificação de dependências por \textit{Software Composition Analysis} e \textit{scanning} de IaC para prevenir falhas de provisionamento.
